\textbf{Phần 1. Trắc nghiệm lựa chọn}
\setcounterpageref{ex}{0}

\begin{ex} %[4X13N72]
Cho dãy số liệu về số cây xanh 4 tổ trồng: $15, 18, 20, 17$. Tổ trồng nhiều cây nhất trồng được
\choice
{$18\text{ cây}$}
{\True $20\text{ cây}$}
{$17\text{ cây}$}
{$15\text{ cây}$}
\loigiai{
$20$ cây là số lớn nhất trong dãy.\\
\textbf{Đáp án của câu hỏi là B}
}
\end{ex}

\begin{ex} %[4X13N72]
Khi tung một đồng xu 1 lần, các khả năng có thể xuất hiện ở mặt trên là
\choice
{Mặt ngửa hoặc mặt nghiêng}
{Chắc chắn mặt ngửa}
{Chắc chắn mặt sấp}
{\True Mặt ngửa hoặc mặt sấp}
\loigiai{
Đồng xu có 2 mặt: ngửa hoặc sấp.\\
\textbf{Đáp án của câu hỏi là D}
}
\end{ex}

\begin{ex} %[4X13H72]
Số trung bình cộng của dãy số liệu: $10, 20, 30, 40$ là
\choice
{\True $25$}
{$20$}
{$30$}
{$35$}
\loigiai{
$(10 + 20 + 30 + 40) : 4 = 25$.\\
\textbf{Đáp án của câu hỏi là A}
}
\end{ex}

\begin{ex} %[4X13H72]
Biểu đồ cột số học sinh yêu thích các môn học: Toán ($24$), Tiếng Việt ($20$), Tiếng Anh ($26$). Môn học được nhiều học sinh yêu thích nhất là
\choice
{Toán}
{Tiếng Việt}
{\True Tiếng Anh}
{Cả ba môn như nhau}
\loigiai{
Tiếng Anh có 26 bạn, cao nhất.\\
\textbf{Đáp án của câu hỏi là C}
}
\end{ex}

\begin{ex} %[4X13H72]
Tung đồng xu 20 lần, có 11 lần xuất hiện mặt ngửa. Số lần xuất hiện mặt sấp là
\choice
{$8$ lần}
{\True $9$ lần}
{$10$ lần}
{$11$ lần}
\loigiai{
$20 - 11 = 9$ lần.\\
\textbf{Đáp án của câu hỏi là B}
}
\end{ex}

\begin{ex} %[4X13H72]
Trong một hộp kín có 10 viên bi đỏ và 5 viên bi xanh. Lấy ngẫu nhiên 1 viên bi. Khả năng lấy được bi đỏ
\choice
{Bé hơn khả năng lấy bi xanh}
{Bằng khả năng lấy bi xanh}
{Không thể xảy ra}
{\True Lớn hơn khả năng lấy bi xanh}
\loigiai{
10 bi đỏ $> 5$ bi xanh nên khả năng lấy bi đỏ cao hơn.\\
\textbf{Đáp án của câu hỏi là D}
}
\end{ex}

\begin{ex} %[4X13V72]
Biểu đồ số sách đọc được trong học kì 2 của 4 bạn: An ($12$), Bình ($15$), Cường ($10$), Dung ($18$). Tổng số sách 4 bạn đã đọc là
\choice
{\True $55$ cuốn}
{$50$ cuốn}
{$60$ cuốn}
{$52$ cuốn}
\loigiai{
$12 + 15 + 10 + 18 = 55$ cuốn.\\
\textbf{Đáp án của câu hỏi là A}
}
\end{ex}

\begin{ex} %[4X13V72]
Gieo một con xúc xắc 30 lần, thấy có 6 lần xuất hiện mặt 1 chấm. Số lần xuất hiện các mặt còn lại là
\choice
{$20$ lần}
{$25$ lần}
{\True $24$ lần}
{$26$ lần}
\loigiai{
$30 - 6 = 24$ lần.\\
\textbf{Đáp án của câu hỏi là C}
}
\end{ex}

\begin{ex} %[4X13V72]
Biểu đồ số điểm 10 của 3 tổ: Tổ 1 ($35$), Tổ 2 ($40$), Tổ 3 ($45$). Trung bình mỗi tổ đạt được số điểm 10 là
\choice
{$38$ điểm}
{\True $40$ điểm}
{$42$ điểm}
{$45$ điểm}
\loigiai{
$(35 + 40 + 45) : 3 = 40$ điểm 10.\\
\textbf{Đáp án của câu hỏi là B}
}
\end{ex}

\begin{ex} %[4X13V72]
Cho dãy số liệu điểm thi: $6, 7, 8, 9, 8, 10, 8, 7, 8$. Số điểm 8 xuất hiện bao nhiêu lần?
\choice
{$3$ lần}
{$5$ lần}
{$2$ lần}
{\True $4$ lần}
\loigiai{
Điểm 8 xuất hiện 4 lần.\\
\textbf{Đáp án của câu hỏi là D}
}
\end{ex}

\begin{ex} %[4X13V72]
Quay vòng quay có 3 màu (Xanh, Đỏ, Vàng) 40 lần: kim chỉ vào màu Đỏ 15 lần, màu Xanh 17 lần. Số lần kim chỉ vào màu Vàng là
\choice
{\True $8$ lần}
{$7$ lần}
{$9$ lần}
{$10$ lần}
\loigiai{
$40 - (15 + 17) = 8$ lần.\\
\textbf{Đáp án của câu hỏi là A}
}
\end{ex}

\begin{ex} %[4X13C72]
Một nhóm 5 học sinh có điểm kiểm tra trung bình là 8 điểm. Nếu thêm điểm của bạn thứ 6 là 9 điểm thì điểm trung bình của cả 6 bạn là bao nhiêu (viết dạng số thập phân)?
\choice
{$8{,}1$}
{$8{,}2$}
{\True $8{,}17$ (hoặc $(40+9)/6 = 49/6$)}
{$8{,}5$}
\loigiai{
Tổng 5 bạn: $8 \times 5 = 40$. Tổng 6 bạn: $40 + 9 = 49$. Trung bình: $49 : 6 \approx 8{,}17$.\\
\textbf{Đáp án của câu hỏi là C}
}
\end{ex}

\textbf{Phần 2. Trắc nghiệm đúng sai}
\setcounterpageref{ex}{0}

\begin{ex} %[4X13H72]
Xét tính đúng sai về các yếu tố thống kê và xác suất:
\choiceTF
{\True Biểu đồ cột giúp trực quan hóa việc so sánh các số liệu.}
{\True Số trung bình cộng nằm giữa giá trị nhỏ nhất và lớn nhất của dãy số liệu.}
{Tung đồng xu 100 lần thì chắc chắn số lần ngửa luôn bằng 50 lần.}
{\True Khi gieo xúc xắc 6 mặt, có 6 khả năng xảy ra đối với số chấm ở mặt trên.}
\loigiai{
\begin{itemize}
    \item Biểu đồ cột trực quan: Đúng.
    \item TBC nằm giữa min và max: Đúng.
    \item Số lần ngửa là ngẫu nhiên, không chắc chắn bằng 50. Sai.
    \item 6 mặt từ 1 đến 6 chấm: Đúng.
\end{itemize}
\begin{itemize}
    \item a) \textbf{Đúng}.
    \item b) \textbf{Đúng}.
    \item c) \textbf{Sai}.
    \item d) \textbf{Đúng}.
\end{itemize}
}
\end{ex}

\begin{ex} %[4X13V72]
Biểu đồ số cây 3 lớp trồng được: 4A ($40\text{ cây}$), 4B ($50\text{ cây}$), 4C ($60\text{ cây}$). Xét tính đúng sai:
\choiceTF
{\True Lớp 4C trồng được nhiều cây nhất.}
{\True Lớp 4A trồng được ít cây nhất.}
{Cả 3 lớp trồng được tất cả $160\text{ cây}$.}
{\True Trung bình mỗi lớp trồng được $50\text{ cây}$.}
\loigiai{
\begin{itemize}
    \item 4C nhiều nhất (60 cây): Đúng.
    \item 4A ít nhất (40 cây): Đúng.
    \item Tổng: $40 + 50 + 60 = 150\text{ cây} \neq 160\text{ cây}$. Sai.
    \item TBC: $150 : 3 = 50\text{ cây}$. Đúng.
\end{itemize}
\begin{itemize}
    \item a) \textbf{Đúng}.
    \item b) \textbf{Đúng}.
    \item c) \textbf{Sai}.
    \item d) \textbf{Đúng}.
\end{itemize}
}
\end{ex}

\textbf{Phần 3. Trả lời ngắn}
\setcounterpageref{ex}{0}

\begin{ex} %[4X13H72]
Tính số trung bình cộng của 3 số: $30, 40, 50$.
\shortans{40}
\loigiai{
\textbf{Đáp số:} $40$\\
$40$.
}
\end{ex}

\begin{ex} %[4X13H72]
Tung đồng xu 40 lần có 22 lần ngửa. Tìm số lần sấp.
\shortans{18}
\loigiai{
\textbf{Đáp số:} $18$\\
$40 - 22 = 18$ lần.
}
\end{ex}

\begin{ex} %[4X13V72]
Cho dãy số liệu: $15, 20, 18, 25, 30$. Tìm số lớn nhất trong dãy.
\shortans{30}
\loigiai{
\textbf{Đáp số:} $30$\\
$30$.
}
\end{ex}

\begin{ex} %[4X13V72]
Biểu đồ số học sinh giỏi 3 lớp: $12, 15, 18$. Tính tổng số học sinh giỏi cả 3 lớp.
\shortans{45}
\loigiai{
\textbf{Đáp số:} $45$\\
$12 + 15 + 18 = 45$.
}
\end{ex}

\begin{ex} %[4X13V72]
Gieo xúc xắc 50 lần có 8 lần ra mặt 6 chấm. Hỏi các mặt còn lại xuất hiện bao nhiêu lần?
\shortans{42}
\loigiai{
\textbf{Đáp số:} $42$\\
$50 - 8 = 42$ lần.
}
\end{ex}

\begin{ex} %[4X13C72]
Bốn bạn có chiều cao trung bình là $135\text{ cm}$. Bạn thứ năm có chiều cao $140\text{ cm}$. Tính chiều cao trung bình của cả 5 bạn theo xăng-ti-mét.
\shortans{136}
\loigiai{
\textbf{Đáp số:} $136$\\
$(135 \times 4 + 140) : 5 = (540 + 140) : 5 = 680 : 5 = 136\text{ cm}$.
}
\end{ex}
